%-------------------------
% Resume in Latex
% Based off of: https://github.com/sb2nov/resume
% License : MIT
%------------------------

\documentclass[letterpaper,10pt]{article}

\usepackage{latexsym}
\usepackage{fontawesome5}
\usepackage[empty]{fullpage}
\usepackage{titlesec}
\usepackage{marvosym}
\usepackage[usenames,dvipsnames]{color}
\usepackage{verbatim}
\usepackage{enumitem}
\usepackage[hidelinks]{hyperref}
\usepackage{fancyhdr}
\usepackage[russian]{babel}
\usepackage{tabularx}
\usepackage[super]{nth}
\usepackage{booktabs}
\usepackage{accsupp}
\usepackage{xifthen}
\usepackage{fontspec}

%----------FONT OPTIONS----------
% Use a standard font with Cyrillic support
\setmainfont{Times New Roman}

\pagestyle{fancy}
\fancyhf{} % clear all header and footer fields
\fancyfoot{}
\fancyhead{}
\renewcommand{\headrulewidth}{0pt}
\renewcommand{\footrulewidth}{0pt}

% Adjust margins
\addtolength{\oddsidemargin}{-0.5in}
\addtolength{\evensidemargin}{-0.5in}
\addtolength{\textwidth}{1in}
\addtolength{\topmargin}{-.5in}
\addtolength{\textheight}{1.0in}

\urlstyle{same}
% defines one's email
% usage: \email{<email>}
\newcommand{\email}[1]{\faEnvelope\ \href{mailto:#1}{#1}}
% defines one's phone
% usage: \phone{<phone>}
\newcommand{\phone}[1]{\faPhone\ {#1}}
% defines one's linkedin
% usage: \linkedin{<linkedin>}
\newcommand{\linkedin}[1]{\href{#1}{\faLinkedin}}
% defines one's GitHub
% usage: \github{<github>}
\newcommand{\github}[1]{\href{#1}{\faGithub}}

% defines one's homepage
% usage: \homepage{<homepage>}
\newcommand{\homepage}[2][]{\faLink\
  \ifthenelse{\isempty{#1}}%
    {\href{#2}{#2}}
    {\href{#2}{#1}}}
\newcommand{\researchgate}[1]{\href{#1}{\faResearchgate}}
\newcommand{\orcid}[1]{\href{#1}{\faOrcid}}

\raggedbottom
\raggedright
\setlength{\tabcolsep}{0in}

% Sections formatting
\titleformat{\section}{
  \vspace{-4pt}\scshape\raggedright\large
}{}{0em}{}[\color{black}\titlerule \vspace{-5pt}]

% Ensure that generate pdf is machine readable/ATS parsable
% \pdfgentounicode=1

%-------------------------
% Custom commands
\newcommand{\link}[2]{\href{#1}{\color{blue}\underline{#2}}}

\newcommand{\resumeItem}[1]{
  \item\small{
    {#1 \vspace{-2pt}}
  }
}

\newcommand{\resumeSubheading}[4]{
  \vspace{-2pt}\item
  \begin{tabular*}{0.97\textwidth}[t]{l@{\extracolsep{\fill}}r}
    \textbf{#1} & #2 \\
    \textit{\small#3} & \textit{\small #4} \\
  \end{tabular*}\vspace{-7pt}
}

\newcommand{\resumeSubSubheading}[2]{
  \item
  \begin{tabular*}{0.97\textwidth}{l@{\extracolsep{\fill}}r}
    \textit{\small#1} & \textit{\small #2} \\
  \end{tabular*}\vspace{-7pt}
}

\newcommand{\resumeProjectHeading}[2]{
  \item
  \begin{tabular*}{0.97\textwidth}{l@{\extracolsep{\fill}}r}
    \small#1 & #2 \\
  \end{tabular*}\vspace{-7pt}
}

\newcommand{\resumeSubItem}[1]{\resumeItem{#1}\vspace{-4pt}}

\renewcommand\labelitemii{$\vcenter{\hbox{\tiny$\bullet$}}$}

\newenvironment{resumeSubHeadingList}{\begin{itemize}[leftmargin=0.15in, label={}]}{\end{itemize}}
\newenvironment{resumeItemList}{\begin{itemize}}{\end{itemize}}
\newenvironment{resumeItemSubList}{\begin{itemize}\setlength\itemsep{0.4em}}{\end{itemize}\vspace{-5pt}}

\newcommand{\awardsTable}[1]{
	\begin{tabularx}{\textwidth}{l @{\extracolsep{\fill}} *{4}{l}}
	\emph{Name}	& \emph{Placing}	& \emph{Scope}	& \emph{Awarder} & \emph{Year}       \\
	\midrule\midrule
	#1
	\end{tabularx}
}
\newcommand{\awardsTableRow}[6]{
	\BeginAccSupp{method=plain, ActualText=11\string\t 21}#1 & #2 & #3 & #4 & #5\EndAccSupp{} \\
}

\fancyfoot[L]{
  \small \LaTeX\ исходный код: \link{https://github.com/vectorsss/cv}{github.com/vectorsss/cv} $|$
  Веб-версия: \link{https://cv.zhaochi.ru/ru}{cv.zhaochi.ru/ru}
}

\renewcommand\footrule{\hrule width\textwidth}

%-------------------------------------------
%%%%%%  RESUME STARTS HERE  %%%%%%%%%%%%%%%%%%%%%%%%%%%%

\begin{document}
\begin{center}
  \begin{flushright}
    \footnotesize{\it Дата компиляции: \today{}}
  \end{flushright}
\end{center}
\begin{center}
  \textbf{\huge \scshape Чжао Чи (Vector)} \\ \vspace{1pt}
  \small \email{vector.zhaochi@gmail.com} $|$ \homepage[https://chizhao.gitlab.io]{https://chizhao.gitlab.io} $|$ \linkedin{https://www.linkedin.com/in/dandanv5} $|$ \github{https://github.com/vectorsss} $|$
  \researchgate{https://www.researchgate.net/profile/Chi-Zhao-4} $|$ \orcid{https://orcid.org/0000-0002-1166-7578}
\end{center}

\section{Образование}
\begin{resumeSubHeadingList}
  \resumeSubheading
    {Санкт-Петербургский государственный университет}{сент. 2021 -- июнь 2025}
    {Аспирантура, прикладная математика и процессы управления}{Санкт-Петербург, Россия}
    \resumeSubSubheading
    {Тема диссертации: <<Моделирование динамики бинарных мнений в социальных сетях сложных конфигураций>>}{Защищена в апр. 2025}

  \resumeSubheading
    {Санкт-Петербургский государственный университет}{сент. 2019 -- июнь 2021}
    {Магистратура, прикладная математика и информатика (средний балл: 4,9/5,0)}{Санкт-Петербург, Россия}

  \resumeSubheading
    {Пекинский технологический институт}{3 -- 28 июля 2017}
    {Летняя школа ACM}{Пекин, Китай}

  \resumeSubheading
    {Яньаньский университет}{2015 -- 2019}
    {Бакалавриат, информатика (средний балл: 3,1/4,0)}{Яньань, Китай}
\end{resumeSubHeadingList}

\section{Опыт работы}
\begin{resumeSubHeadingList}
  \resumeSubheading
    {Санкт-Петербургский государственный университет}{дек. 2025 -- н.в.}
    {Ассистент}{Санкт-Петербург, Россия}
  \resumeSubheading
    {Huawei Inc.}{сент. 2021 -- дек. 2025}
    {Инженер-исследователь}{Санкт-Петербург, Россия}
  \resumeSubheading
    {Wisedu Inc.}{сент. 2016 -- май 2019}
    {Инженер-программист}{Яньань, Китай}

\end{resumeSubHeadingList}
\section{Проектный опыт}
\begin{resumeSubHeadingList}
  \resumeSubheading
    {* Управление данными моделей / Отбор ценных обучающих данных}{2025}
    {Руководитель проекта}{Санкт-Петербург, Россия}
    \resumeSubSubheading{Повышение точности моделей через управление данными и оптимизацию отбора признаков.}{}
    \resumeSubSubheading{Выборка 60 млн записей из 1,8 млрд с помощью оригинальных алгоритмов.}{}
    \resumeSubSubheading{Повышение точности модели на 22\% через алгоритмы выборки и обнаружения аномалий.}{}
    \resumeSubSubheading{Ускорение обучения на 10\% за счёт удаления 20\% избыточных признаков.}{}
  \resumeSubheading
    {* Внутренний алгоритм колоночного хранения}{2023--2024}
    {Инженер-исследователь}{Санкт-Петербург, Россия}
    \resumeSubSubheading{Высокопроизводительный, гибкий алгоритм сжатия без потерь для колоночного хранения}{}
    \resumeSubSubheading{Алгоритм будет коммерциализирован для хранения данных базовых станций.}{}
    \resumeSubSubheading{LTE (4G) --- 30\% сжатие, NR (5G) --- 40\% сжатие, без потери производительности.}{}
  \resumeSubheading
    {Грант Российского научного фонда \link{https://rscf.ru/en/project/22-21-00346/}{№ 22-21-00346}}{2023}
    {Исследователь}{Санкт-Петербург, Россия}
    \resumeSubSubheading{Моделирование двухслойных моделей динамики мнений: выражаемые и скрытые мнения.}{}
  \resumeSubheading
    {* Алгоритм эластичного расширения для распределённых файловых хранилищ}{2023}
    {Инженер-исследователь}{Санкт-Петербург, Россия}
    \resumeSubSubheading{Бесшовное расширение HOFS через HRW-хеширование с минимальным перемещением данных.}{}
    \resumeSubSubheading{Повышена надёжность: непрерывный доступ к данным даже после потери хеш-кольца.}{}
  \resumeSubheading
    {* Сжатие данных пакетов/маршрутизаторов}{2023}
    {Инженер-исследователь}{Санкт-Петербург, Россия}
    \resumeSubSubheading{Сжатие 85\% без потерь (экономия дискового пространства в 7 раз).}{}
  \resumeSubheading
    {* Сжатие беспроводных данных}{2021--2022}
    {Инженер-исследователь}{Санкт-Петербург, Россия}
    \resumeSubSubheading{Достигнута степень сжатия 96\% с помощью сжатия с потерями.}{}
  \resumeSubheading
    {* Оптимизация запросов SparkSQL}{2021}
    {Инженер-исследователь}{Санкт-Петербург, Россия}
    \resumeSubSubheading{Ускорение запросов на 50\% через проталкивание предикатов и оптимизацию структур данных.}{}
\end{resumeSubHeadingList}

\newpage
%-----------PROJECTS-----------
\section{Патенты / Свидетельства о регистрации}
\begin{resumeSubHeadingList}
  \resumeSubheading
    {Программа для моделирования динамики распространения бинарных мнений в двухслойных сетях}{2023}
    {Свидетельство о государственной регистрации программы для ЭВМ (Россия)}{{\link{https://new.fips.ru/registers-doc-view/fips_servlet?DB=EVM&DocNumber=2023661532&TypeFile=html}{№ 2023661532}}}
  \resumeSubheading
    {Система классификации онлайн-новостей на основе свёрточной нейронной сети}{2018}
    {Свидетельство о регистрации программного обеспечения (Китай)}{№ 2831192}
\end{resumeSubHeadingList}

\section{Проекты с открытым исходным кодом}
\begin{resumeSubHeadingList}

  \resumeProjectHeading{
    \link{https://github.com/vectorsss/shapG}{\textbf{ShapG}} $|$
    \emph{python $\cdot$ алгоритм важности признаков $\cdot$ меры
          центральности}}{}
  \begin{resumeItemList}
    \resumeItem{Python-пакет для алгоритмов важности признаков на основе
                значений Шепли.}
    \resumeItem{Пакет также может использоваться для вычисления центральности
      в графах.}
  \end{resumeItemList}
  \resumeProjectHeading{
    \link{https://github.com/vectorsss/news_classification}{\textbf{News
        Classification}} $|$
    \emph{python $\cdot$ NLP $\cdot$ CNN $\cdot$ классификация текстов}}{}
  \begin{resumeItemList}
    \resumeItem{Система классификации онлайн-новостей на основе свёрточной нейронной
      сети.}
    \resumeItem{Проект получил 3-е место на провинциальном этапе Всекитайского
      конкурса компьютерного проектирования 2018 года.}
  \end{resumeItemList}
\end{resumeSubHeadingList}

\section{Навыки}
\begin{itemize}[leftmargin=0.15in, label={}]
  \small{\item{
        \textbf{Языки программирования}{: Python $\cdot$ C/C++ $\cdot$ Go $\cdot$ Rust $\cdot$ Matlab/Octave $\cdot$ Julia $\cdot$ R $\cdot$ SQL $\cdot$ \LaTeX} \\
        \textbf{Библиотеки МО}{: Tensorflow $\cdot$  PyTorch $\cdot$ Keras $\cdot$  Scikit-Learn} \\
        \textbf{Инструменты разработки}{: Git $\cdot$ Docker $\cdot$  Google Cloud Platform} \\
        \textbf{ОС}{: MacOS $\cdot$ Arch Linux $\cdot$ Windows} \\
        \textbf{Языки}{: китайский (родной), английский (свободное владение), русский (свободное владение)}
        }}
\end{itemize}

\section{Научные интересы}
\begin{resumeSubHeadingList}
  \resumeProjectHeading{
    Алгоритмы на графах $\cdot$ Меры центральности $\cdot$ Машинное обучение
    $\cdot$ Объяснимый ИИ $\cdot$ Теория вероятностей $\cdot$ Статистика}{}
  \resumeProjectHeading{
    Сжатие данных $\cdot$ Теория кодирования $\cdot$ Анализ временных рядов $\cdot$
    Оптимизация $\cdot$ Стохастическое моделирование $\cdot$ Случайные процессы}{}
\end{resumeSubHeadingList}
\section{Преподавание}
\begin{resumeSubHeadingList}
  \resumeSubheading{Optimization, Forecasting and Artificial Intelligence in
                    Industry}{Санкт-Петербург, Россия}{Преподаватель}{2026}
  \resumeSubheading{Deep learning}{Санкт-Петербург,
                    Россия}{Преподаватель}{2026}
  \resumeSubheading{Advance Economics}{Санкт-Петербург,
                    Россия}{Преподаватель}{2025}
  \resumeSubheading{Applied Statistics in R}{Санкт-Петербург, Россия}{Ассистент
                    преподавателя}{2024}
  \resumeSubheading{Statistical Decisions and Econometrics}{Санкт-Петербург,
                    Россия}{Ассистент преподавателя}{2022}
\end{resumeSubHeadingList}
\section{Научная деятельность}
\begin{resumeSubHeadingList}
  \resumeProjectHeading{{\bf Рецензент журналов и конференций}}{}
  \begin{resumeItemList}
    \resumeItem{Engineering Applications of Artificial Intelligence (EAAI)}
    \resumeItem{Group Decision and Negotiation (GDN)}
    \resumeItem{Dynamic Games and Applications (DGAA)}
    \resumeItem{Humanities and Social Sciences Communications (HSSCOMMS)}
    \resumeItem{International Conference On Computational Optimization (ICOMP)}
  \end{resumeItemList}
\end{resumeSubHeadingList}
\section{Избранные публикации}
\begin{enumerate}
  \fontsize{10}{10.5}\selectfont
  \item \textbf{Zhao C.}, Liu J., Parilina E. M. Strategic inputs: feature selection from game-theoretic perspective. // {\it arXiv preprint arXiv:2510.24982.} -- 2025.
  \item \textbf{Zhao C.}, Liu J., Parilina E. M. The Shapley Value Contribution to Explainable Artificial Intelligence: A Comprehensive Survey // {\it Dynamic Games and Applications.} -- 2025. -- Oct. -- P. 1-38. {\bf (Q2)}
  \item \textbf{Zhao C.}, Liu J., Parilina E. M. Complete-to-Sparse: A Novel Graph Construction Strategy for Efficient ShapG // {\it Mathematical Optimization Theory and Operations Research}. -- Cham : Springer Nature Switzerland. -- 2025. -- P. 180–194.
  \item \textbf{Zhao C.}, Liu J., Parilina E. M. ShapG: new feature importance method based on the Shapley value // {\it Engineering Applications of Artificial Intelligence.} -- 2025. -- May. -- Vol. 148, 110409. {\bf (Q1, IF: 8.0)}
        \newpage
  \item \textbf{Zhao C.}, Parilina E. M. Centrality measures and opinion dynamics in two-layer networks with replica nodes // {\it Computers and Operations Research.} -- 2026. -- Jan. -- Vol. 185, 107245. {\bf (Q1, IF: 4.3)}
  \item \textbf{Zhao C.}, Parilina E. M. Analysis of consensus time and winning rate in two-layer networks with hypocrisy of different structures // {\it Vestnik of Saint Petersburg University. Applied Mathematics. Computer Science. Control Processes.} -- 2024. -- Vol. 20, no. 2. -- P. 170-192.
  \item \textbf{Zhao C.}, Parilina E. M. Opinion Dynamics in Two-Layer Networks with Hypocrisy // {\it Journal of the Operations Research Society of China}. -- 2024. -- Mar. -- Vol. 12, no. 1. -- P. 109-132. {\bf (Q2)}
  \item \textbf{Zhao C.}, Parilina E. M. Network Structure Properties and Opinion Dynamics in Two-Layer Networks with Hypocrisy // {\it Mathematical Optimization Theory and Operations Research}. -- Cham : Springer Nature Switzerland. -- 2024. -- P. 300-314.
  \item \textbf{Zhao C.}, Parilina E. M. Consensus time and winning rate based on simulations in two-layer networks with hypocrisy // {\it 2023 7th Scientific School Dynamics of Complex Networks and their Applications (DCNA).} -- 2023. -- P. 68-71.
\end{enumerate}

\section{Конференции}
\begin{resumeSubHeadingList}
  \resumeSubheading{International Conference On Computational Optimization
                    (ICOMP 2025)}{Abu Dhabi, UAE}{Poster Presentation}{Oct. 17
                    - 19, 2025}
  \resumeSubheading{Mathematical Optimization Theory and Operations Research
                    (MOTOR 2025)}{Novosibirsk, Russia}{Oral Presentation}{July.
                    07 - 11, 2025}
  \resumeSubheading{Game Theory and Management (GTM 2025)}{Saint-Petersburg,
                    Russia}{Oral Presentation}{July. 2 - 4, 2025}
  \resumeSubheading{\nth{14} International Society of Dynamic Games (ISDG) Workshop}{Yerevan, Armenia}{Oral Presentation}{June. 11 - 13, 2025}
  \resumeSubheading{International Conference On Computational Optimization
                    (ICOMP 2024)}{Innopolis, Russia}{Visitor}{Oct. 10 - 12,
                    2024}
  \resumeSubheading{Mathematical Optimization Theory and Operations Research
                    (MOTOR 2024)}{Omsk, Russia}{Oral Presentation}{June. 30 -
                    July. 06, 2024}
  \resumeSubheading{Game Theory and Management (GTM 2024)}{Saint-Petersburg,
                    Russia}{Oral Presentation}{June. 26 - 28, 2024}
  \resumeSubheading{Dynamics of Complex Networks and their Applications (DCNA
                    2023)}{Kaliningrad, Russia}{Poster Presentation}{Sep. 18 -
                    20, 2023}
  \resumeSubheading{Game Theory and Management (GTM 2023)}{Saint-Petersburg,
                    Russia}{Oral Presentation}{June. 28 - 30, 2023}
  \resumeSubheading{Control Processes and Stability 2022}{Saint-Petersburg,
                    Russia}{Oral Presentation}{Apr. 4 - 8, 2022}
  \resumeSubheading{Control Processes and Stability 2021}{Saint-Petersburg,
                    Russia}{Oral Presentation}{Apr. 5 - 9, 2021}
    % \resumeSubheading{The Computing Conference 2017}{Hangzhou,
    % China}{Visitor}{Oct. 11 - 14, 2017} \resumeSubheading{Yiban Developer
    % Conference 2017}{Shanghai, China}{Developer}{Aug. 2017}
    % \resumeSubheading{Language \& Intelligence Summit 2017}{Beijing,
    % China}{Visitor}{July 23, 2017}
\end{resumeSubHeadingList}

\section{Награды и достижения}
\setlength{\tabcolsep}{6.0pt}
\renewcommand{\arraystretch}{1.1}
\fontsize{8.5}{11}\selectfont
\awardsTable{
  \awardsTableRow{Outstanding contributor}{-}{Huawei Inc.}{Huawei Chebyshev
                  Research Center ADN AI Lab}{Dec. 2025} \\
  \awardsTableRow{AI4OPT Competition Best Paper}{\nth{1}}{ICOMP 2025}{ICOMP
                  2025 Conference Organizing Committee}{Oct. 2025} \\
  \awardsTableRow{General Development Star}{-}{Huawei Inc.}{General development
                  department}{Oct. 2024} \\
  \awardsTableRow{General Development Star}{-}{Huawei Inc.}{General development
                  department}{Jun. 2023} \\
  \awardsTableRow{General Development Star}{-}{Huawei Inc.}{General development
                  department}{Oct. 2022} \\
  \awardsTableRow{President's Award (Team)}{-}{Huawei Inc.}{MAE-M
                  department}{Dec. 2021} \\
  \awardsTableRow{Diploma with distinction}{-}{University}{Saint Petersburg
                  State University}{Jun. 2021} \\
  \awardsTableRow{Excellent Graduation Thesis (Design)}{-}{University}{Yanan
                  University}{Jun. 2019} \\
  \awardsTableRow{Excellent Graduate}{-}{University}{Yanan University}{Jun.
                  2019} \\
  \awardsTableRow{Merit Student Scholarship}{-}{University}{Yanan
                  University}{Dec. 2018} \\
  \awardsTableRow{China Software Cup Competition}{\nth{3}}{China}{China
                  Software Cup Organizing Committee}{Oct. 2018} \\
  \awardsTableRow{Computer Design Competition}{\nth{3}}{Northwest
                  China}{Northwest University (China)}{May. 2018} \\
  \awardsTableRow{Mathematical Contest in Modeling}{\nth{2}}{Shaanxi}{China
                  Society for Industrial and Applied Mathematics}{Dec. 2017} \\
  \awardsTableRow{Mathematics Competition}{\nth{3}}{Shaanxi}{Chinese
                  Mathematical Society}{Nov. 2016} \\ \hline }

\end{document}
